% ============================================================
%  附录 D：从这场演讲到 Kapustin--Witten
% ============================================================
\clearpage
\markboth{附录 D　从这场演讲到 Kapustin--Witten}{附录 D　从这场演讲到 Kapustin--Witten}
\section*{附录 D　从这场演讲到 Kapustin--Witten：一份技术路线图}
\addcontentsline{toc}{section}{附录 D　从这场演讲到 Kapustin--Witten：一份技术路线图}
\setcounter{equation}{0}\setcounter{figure}{0}\setcounter{table}{0}
\renewcommand{\theequation}{D.\arabic{equation}}
\renewcommand{\thefigure}{D.\arabic{figure}}
\renewcommand{\thetable}{D.\arabic{table}}

\begin{quote}
\small
本附录不属于演讲原文，也不属于逐段批注，而是一份补充地图：把第 6--9 章的线索按
Kapustin--Witten（2006）的逻辑重排一遍，并标出每一步在数学上对应什么。读者若只想
知道“结论是什么”，读 \S D.1 与 \S D.5 即可；想核对公式，再读中间两节。
\end{quote}

\subsection*{D.1　为什么对偶必须写成 $\SL(2,\ZZ)$}

第 3 章里 Montonen--Olive 的对偶写作 $e^{\vee\,2}=4\pi/e^{2}$。把这个公式放进
$\cN=4$ 超对称规范理论的语境，它会被一个更整齐的对象取代。定义复耦合
\begin{equation}
  \tau=\frac{\theta}{2\pi}+\frac{4\pi\ii}{g^{2}},
  \label{eq:tau}
\end{equation}
其中 $\theta$ 是 Yang--Mills 的 $\theta$ 角，$g$ 是规范耦合。$\cN=4$ 理论有两个基本
的“对偶变换”：
\begin{itemize}
  \item $T:\tau\mapsto\tau+1$，即 $\theta\mapsto\theta+2\pi$。这是 $\theta$ 的周期性，
        不改变规范群，也不交换电与磁；
  \item $S:\tau\mapsto-\dfrac{1}{n_g\tau}$，其中 $n_g$ 是“单边性”因子：单连通
        （单边）群 $n_g=1$，$F_4$、$\Sp(k)$、$\mathrm{Spin}(2k+1)$ 取 $n_g=2$，$G_2$ 取
        $n_g=3$。$S$ 同时把规范群 $G$ 换成对偶群 $\gvee$，把电荷换成磁荷。
\end{itemize}
$S$ 与 $T$ 生成整个 $\SL(2,\ZZ)$。在超荷上这个作用不是线性的，而是一个手征旋转，
相因子为 $\bigl(|c\tau+d|/(c\tau+d)\bigr)^{1/2}$，所以真正作用在希尔伯特空间上的是
$\SL(2,\ZZ)$ 的二重覆盖。1979 年时这些结构还完全看不清：$\cN=4$ 理论本身尚未被
系统地写下，$\SL(2,\ZZ)$ 对偶只是一个关于质量谱的猜想。

在 $\cN=2$ 的情形，同样的 $\SL(2,\ZZ)$ 作用在电荷格 $H^{1}(T^{2},\ZZ)$ 上，$(m,n)$
是电磁荷；BPS 态的质量由中心荷给出
\begin{equation}
  Z=\sqrt{\frac{2}{\Im\tau}}\;\tilde\varphi\cdot(\tilde n+\tau\tilde m),
  \qquad M=|Z|,
  \label{eq:bps}
\end{equation}
其中 $\tilde\varphi$ 是标量场的真空期望值。式 \eqref{eq:bps} 是第六章 Sen 论证的
出发点：对偶预言某些 $(m,n)$ 的束缚态存在，而它们的质量由 \eqref{eq:bps} 唯一确定。

\subsection*{D.2　't Hooft 算子：对偶群的表示}

第 8 章已经说明，'t Hooft 算子的“型”由同态
\[
  \rho:U(1)\longrightarrow G
\]
的共轭类给出（等价地，由 $G$ 的一个余权，模去 Weyl 群的作用）。Kapustin 的关键
观察是：这个型恰好等价于对偶群 $\gvee$ 的一个不可约表示。理由是一行根数据的对偶：
\[
  \Hom\bigl(U(1),G\bigr)_{\text{共轭}}\;\cong\;\Hom\bigl(T,\CC^{\times}\bigr)/W
  \;\cong\;\text{$\gvee$ 的支配权},
\]
而后者正是 $\gvee$ 的不可约表示的分类数据（最高权）。于是可以把 't Hooft 算子记作
$T(\gvee R,S)$，其中 $\gvee R$ 是 $\gvee$ 的不可约表示、$S$ 是它所在的曲线。

\begin{table}[htbp]
\centering\small
\caption{线算子的三种语言。左列是物理，中列是群论，右列是几何 Langlands。}
\label{tab:three-languages}
\renewcommand{\arraystretch}{1.35}
\begin{tabular}{@{}p{3.5cm}p{4.6cm}p{5.2cm}@{}}
\toprule
物理对象 & 群论数据 & 几何侧的对象 \\
\midrule
Wilson 算子 $W(R,S)$ & $G$ 的不可约表示 $R$ & $\Bun_{G}$ 上的“电”Hecke 算子 \\
't Hooft 算子 $T(\gvee R,S)$ & $\gvee$ 的不可约表示 $\gvee R$ & Hecke 修改，型为 $\gvee R$ 的最高权 \\
电磁对偶 & $G\leftrightarrow\gvee$ & 镜像对称 \\
\bottomrule
\end{tabular}
\end{table}

表 \ref{tab:three-languages} 就是整场演讲的骨架：三种语言说的是同一件事，而
“Wilson 算子 $\leftrightarrow$ 't Hooft 算子”的对偶，在几何侧就是
“Hecke 算子 $\leftrightarrow$ Hecke 修改”的对应。

\subsection*{D.3　从线算子到 Hecke 算子}

设 $C$ 是紧 Riemann 面，$G$ 是复约化群。$\Bun_{G}$ 是 $C$ 上 $G$-丛的模叠。在点
$z\in C$ 处，Hecke 对应是
\[
  \cH_{G,z}\ \xrightarrow{\ p_{1}\ }\ \Bun_{G},
  \qquad
  \cH_{G,z}\ \xrightarrow{\ p_{2}\ }\ \Bun_{G},
\]
其中 $\cH_{G,z}$ 参数化“在 $z$ 处相差一个 Hecke 修改”的丛对 $(\cE,\cE')$，两个投影
分别取出 $\cE$ 与 $\cE'$。给定
$\gvee$ 的表示 $V$，Hecke 函子是
\[
  \mathbb{H}_{V,z}:\ \cF\ \longmapsto\ p_{2*}\bigl(p_{1}^{*}\cF\otimes V\bigr),
\]
即“拉回—张量—前推”。这正是第 7 章几何 Hecke 变换的层论版本：$V$ 的最高权决定
修改的型，$\Bun_{G}$ 上的层被送到另一层。

Kapustin--Witten 的第 9 节证明：在扭化后的 $\cN=4$ 理论中，沿 $S$ 插入一个
't Hooft 算子，在二维边界理论里恰好实现 $\mathbb{H}_{V,z}$ 的作用。也就是说，
物理里的“ disorder 算子”与代数几何里的“Hecke 函子”是同一个东西。这一步把
第 8 章的线算子与第 7 章的 Hecke 变换接上了。

\subsection*{D.4　扭化、Hitchin 模空间与镜像对称}

把 $\cN=4$ 超杨--米尔斯理论放在 $\RR\times\Sigma\times C$ 上（$\Sigma$ 是二维
“时间—空间”面，$C$ 是紧 Riemann 面），并做一个拓扑扭化。扭化后，理论在 $\Sigma$
方向上只看得到一个二维 sigma 模型，其靶空间是 Hitchin 模空间
\[
  \cM_{H}(G,C)=\Bigl\{(\cE,\varphi)\ \Bigm|\ \varphi\in H^{0}\bigl(C,\ad\cE\otimes K_{C}\bigr)\Bigr\}
  \Big/\!\!\!\Big/\,G,
\]
即“Higgs 丛”的模空间。它是超 Kähler 的，复维数为
\[
  \dim_{\CC}\cM_{H}(G,C)=(2g-2)\dim G .
\]
三个 Kähler 形式对应三个矩映射 $\mu_{I},\mu_{J},\mu_{K}$，超 Kähler 商
$\cM_{H}=\mu^{-1}(0)/G$ 就是第 5 章单极子模空间的“紧 Riemann 面版本”。

在这个图像里：
\begin{itemize}
  \item 线算子（Wilson 与 't Hooft）对应 $\cM_{H}$ 上的膜；
  \item $S$-对偶把 $G$ 换成 $\gvee$，同时把 $\cM_{H}(G,C)$ 与
        $\cM_{H}(\gvee,C)$ 互换，并把 A-模型换成 B-模型——这就是镜像对称；
  \item 几何 Langlands 的“特征层 $\leftrightarrow$ 局部系统”对应
        “A-膜 $\leftrightarrow$ B-膜”。特征层是 $\Bun_{G}$ 上的 D-模，局部系统是
        $C$ 上 $\gvee$ 的局部系统；在物理侧，前者是 't Hooft 算子作用后的边界态，
        后者是 Higgs 场的谱曲线所定义的“磁”对象。
\end{itemize}

\subsection*{D.5　还差的一步，以及 2024 年}

上面这套图像有一个技术缺口：物理给出的膜住在 $\cM_{H}(G,C)$ 上，而几何 Langlands
要的 D-模住在 $\Bun_{G}$ 上。两者不是同一个空间。Kapustin--Witten 只在“非分歧”
情形、并且在某些极限下把两边对上；完整的对应需要额外的构造。Witten 2009 年的
《Geometric Langlands from six dimensions》用六维 $(2,0)$ 理论给出了一条更直接的
路线：把 $C$ 上的局部系统实现为六维理论中的表面算子，Hecke 特征层则从边界条件
读出。

2024 年，Gaitsgory 与十余位合作者发表五篇论文，证明了几何 Langlands 猜想（非分歧
情形）。这个证明是纯代数几何的，与上面的物理图像并非逐条对应，但两者共享同一个
中心对象：仿射 Grassmannian 上的几何 Satake 等价，以及 $\Bun_{G}$ 上的 Hecke
特征层。对本书的读者而言，值得记住的是：Witten 在 2021 年演讲结尾说“我试着描述
这个故事如何与数学、物理中的其他发展并行展开”——三年后，这条并行线在数学一侧
到达了一个里程碑。物理一侧的解释（为什么偏偏是这套镜像对称）至今仍是开放的。
